% --- INICIO DE PLANTILLA PERSONALIZADA ---
\documentclass[oneside,11pt]{Latex/Classes/PhDthesisPSnPDF}

% --- Paquetes Originales ---
\usepackage{blindtext}
\usepackage{actuarialangle} 
\usepackage{tabularx}
\usepackage{float}
\usepackage{multirow, array}
\usepackage[usenames,dvipsnames,svgnames,table]{xcolor}
\usepackage{longtable}
\usepackage{latexsym,amsmath,amssymb,amsfonts}
\usepackage{enumitem}
\usepackage{xparse}
\usepackage{booktabs}
\usepackage[round, sort, numbers]{natbib}  
\usepackage{adjustbox}
\usepackage[utf8]{inputenc}
\usepackage{caption}
\usepackage{graphicx}
\usepackage{hyperref}
\usepackage{cleveref}

% Definición de colores
\definecolor{miblue}{RGB}{68, 114, 196}

% --- CAMBIO CRÍTICO 1: Usar \input para macros en el preámbulo ---
% \include da problemas antes del begin document. \input es seguro.
\input{Latex/Macros/MacroFile1}

% --- Definiciones de seguridad (por si la clase falla al cargar alguna variable) ---
\providecommand{\facultad}[1]{\def\lafacultad{#1}}
\providecommand{\escudofacultad}[1]{\def\elescudo{#1}}
\providecommand{\degree}[1]{\def\elgrado{#1}}
\providecommand{\director}[1]{\def\eldirector{#1}}
\providecommand{\lugar}[1]{\def\ellugar{#1}}
\providecommand{\degreedate}[1]{\def\lafecha{#1}}

% --- Configuración para bloques de código de R (Pandoc) ---

%%%%%%%%%%%%%%%%%%%%%%%%%%%%%%%%%%%%%%%%%%%%%%%%%%%%%%%%%%%%%%%%%%%%%%%%%%%%%%%%
%                         DATOS (Conectados a RMarkdown)                       %
%%%%%%%%%%%%%%%%%%%%%%%%%%%%%%%%%%%%%%%%%%%%%%%%%%%%%%%%%%%%%%%%%%%%%%%%%%%%%%%%
\title{Analisis del Mercado de Videojuegos}
\author{Carlos Herrera \\ Juan Hernandez}

% Lógica para Facultad: Si la defines en el YAML la usa, si no, usa la default

\facultad{Facultad de Ciencias Económicas y Sociales \\ Escuela de
Estadistica y Ciencias Actuariales}
  \escudofacultad{Latex/Classes/Escudos/faces}

\degree{Computación I}
\director{Prof.~Jesus Ochoa \\ Prof.~Oliver Triveño}
\degreedate{Abril 2026} 
\lugar{Caracas}

\portadatrue 

% Metadatos PDF
\keywords{}
\subject{}

% --- CORRECCIÓN PARA PANDOC NUEVO (Imágenes) ---
\makeatletter
\providecommand{\pandocbounded}[1]{#1}
\makeatother


%%%%%%%%%%%%%%%%%%%%%%%%%%%%%%%%%%%%%%%%%%%%%%%%%%%%%
%                   DOCUMENTO                       %
%%%%%%%%%%%%%%%%%%%%%%%%%%%%%%%%%%%%%%%%%%%%%%%%%%%%%
\begin{document}

\maketitle

%%%%%%%%%%%%%%%%%%%%%%%%%%%%%%%%%%%%%%%%%%%%%%%%%%%%%
%                  PRÓLOGO                          %
%%%%%%%%%%%%%%%%%%%%%%%%%%%%%%%%%%%%%%%%%%%%%%%%%%%%%
\frontmatter

% Puedes agregar dedicatorias desde el YAML si quieres, o dejarlas fijas aquí

%%%%%%%%%%%%%%%%%%%%%%%%%%%%%%%%%%%%%%%%%%%%%%%%%%%%%
%                   ÍNDICES                         %
%%%%%%%%%%%%%%%%%%%%%%%%%%%%%%%%%%%%%%%%%%%%%%%%%%%%%
\setcounter{secnumdepth}{3}
\setcounter{tocdepth}{3}

\tableofcontents



\cleardoublepage

  \cleardoublepage       % Empieza en página derecha (estándar en tesis)
  \chapter*{Resumen}     % Crea el título "Resumen" sin numeración (Capítulo *)
  \addcontentsline{toc}{chapter}{Resumen} % (Opcional) Lo añade al índice
  
  En el presente informe se realiza un análisis exploratorio de los
datos de precios para el año 2023 de una selección de carros. Partiendo
de variables como marca, modelo, año, tipo de transmisión, millaje,
entre otras; se examinan cuales los factores que afectan su precio,
aplicando técnicas estadísticas descriptivas y visualización de datos
con R.             % Aquí se pega el texto que escribiste en el YAML
  
  \cleardoublepage

%%%%%%%%%%%%%%%%%%%%%%%%%%%%%%%%%%%%%%%%%%%%%%%%%%%%%
%                   CONTENIDO (El Corazón)          %
%%%%%%%%%%%%%%%%%%%%%%%%%%%%%%%%%%%%%%%%%%%%%%%%%%%%%
\mainmatter
\def\baselinestretch{1.5}

% --- CAMBIO CRÍTICO 2: Aquí RMarkdown inyecta todo tu texto ---
\chapter*{Introducción}

A lo largo de la vida de una persona, pocas decisiones económicas son
tan importantes como la compra de una casa o un carro. El en el caso de
un carro, el mercado automotriz es altamente dinámico y se caracteriza
por la interacción de múltiples variables a la hora de determinar el
valor de un vehículo. A diferencia de otros bienes, como un inmueble, un
carro es un activo que suele sufrir de una rápida depreciación; cuyo
valor, además de verse afectado por el uso, también podria verse
afectado por factores como avances tecnológicos, prestigio de la marca,
modelo especifico, entre otros. \newline \newline Hoy en día, factores
como el sistema de propulsión; ya sean los tradicionales, como gasolina
o diésel, hasta las nuevas tecnologías hibridas y eléctricas, agregan
nuevas dimensiones a la hora de determinar el valor de un carro.
Entender como estas características técnicas, junto con el uso y la
antigüedad del carro afectan su precio se vuelve entonces de suma
importancia para aquellos consumidores o empresas que quieran maximizar
u optimizar su inversión. \newline \newline Es por ello que en el
presente trabajo se busca aplicar la estadística descriptiva para
explorar y caracterizar el mercado de vehículos, tratando de identificar
patrones que nos permitan tener una visión más clara y objetiva sobre
cuáles son los factores que afectan el valor de un carro

\section*{Estructura del Informe}

Este documento técnico se encuentra estructurado en cinco capítulos
fundamentales:

\begin{itemize}
    
    \item En el \textbf{Capítulo I}, se plantea la problemática de la volatilidad de precios en el mercado automotriz de segunda mano y el impacto de la depreciación técnica, formulando las preguntas de investigación y delimitando el alcance del estudio basado en las marcas y modelos analizados.
    
    \item El \textbf{Capítulo II} establece el Marco Teórico, definiendo los conceptos de estadística descriptiva fundamentales como las medidas de tendencia central, los cuartiles para la detección de valores atípicos y la correlación de Pearson, necesarios para interpretar la relación entre el kilometraje y el valor de mercado.

    \item El \textbf{Capítulo III} describe el Marco Metodológico, detallando el proceso de limpieza y preparación de datos (ETL) realizado en R, el tratamiento de valores ausentes (especialmente en vehículos eléctricos) y los criterios de segmentación por tipo de combustible y transmisión.

    \item El \textbf{Capítulo IV} presenta el Análisis de Resultados, donde se exhiben las tablas de frecuencia y los gráficos (histogramas y diagramas de caja) generados dinámicamente, discutiendo los hallazgos sobre la distribución de precios por marca y la eficiencia de los motores según su tamaño.

    \item Finalmente, en el \textbf{Capítulo V}, se exponen las Conclusiones y Recomendaciones estratégicas derivadas del análisis cuantitativo, orientadas a la valoración precisa de activos y a la optimización de decisiones de compra y venta basadas en evidencia estadística.
\end{itemize}

Este reporte no solo pretende ofrecer una radiografía del estado actual
del mercado de vehículos usados, sino también establecer una metodología
reproducible para el análisis de activos físicos sujetos a depreciación
temporal y tecnológica.

\chapter{Planteamiento del Problema}

En el mercado automotriz actual, el valor de un vehículo no es una cifra
estática, sino el resultado de una interacción multivariable donde
convergen la obsolescencia tecnológica, el desgaste físico y las
preferencias del consumidor. Para un analista o un comprador, la
imprevisibilidad de los precios de reventa representa un riesgo
financiero significativo. \newline \newline El conjunto de datos bajo
estudio presenta una muestra heterogénea que incluye desde vehículos de
combustión tradicional hasta modelos eléctricos de alta gama. Esta
diversidad genera una alta dispersión en los precios, dificultando la
identificación de un ``valor justo''. El problema radica en la falta de
una caracterización estadística clara que permita entender cómo factores
como el kilometraje acumulado, el año de fabricación y el tipo de
combustible (especialmente ante el auge de los motores híbridos y
eléctricos) impactan directamente en la valoración del activo. \newline
\newline Sin un análisis descriptivo robusto, las decisiones de compra y
venta se basan en percepciones subjetivas y no en el comportamiento real
de los datos, lo que justifica la necesidad de este estudio.

Para orientar el análisis, se plantean las siguientes interrogantes:

\textbf{1.} ¿Cuál es el perfil de distribución de los precios de los
vehículos y qué tan alta es su variabilidad interna?

\textbf{2.} ¿Existe una relación inversamente proporcional clara entre
el kilometraje y el precio en todas las marcas por igual?

\textbf{3.} ¿Cómo se comportan los precios de los vehículos según su
tipo de energía (Gasolina vs.~Eléctrico/Híbrido)?

\section{Justificación}

La presente investigación se justifica desde tres dimensiones
fundamentales: la económica, la técnica y la académica. \newline
\newline Desde la dimensión económica, la adquisición de un vehículo
representa una de las inversiones más significativas para el patrimonio
personal o empresarial. No obstante, el mercado de automóviles de
segunda mano se caracteriza por una alta asimetría de información, donde
el precio suele estar sujeto a la especulación o a una valoración
empírica del estado del bien. Este estudio proporciona una base objetiva
que permite comprender la estructura de costos y la tasa de
depreciación, facilitando una toma de decisiones informada que minimice
el riesgo financiero. \newline \newline Desde la dimensión técnica, la
industria automotriz atraviesa una transición energética sin
precedentes. La coexistencia de motores de combustión interna con
tecnologías híbridas y eléctricas plantea un reto para los métodos de
valoración tradicionales. La aplicación de herramientas de
\textbf{Estadística Descriptiva} en este contexto permite caracterizar
estas nuevas categorías y determinar si su comportamiento de mercado
sigue los patrones históricos o si, por el contrario, requieren de
marcos de análisis diferenciados debido a su naturaleza tecnológica.
\newline \newline Desde la dimensión académica, este trabajo se inserta
en la formación de la Escuela de Estadística y Ciencias Actuariales,
demostrando la capacidad del lenguaje de programación R para el manejo
de grandes volúmenes de datos y la generación de reportes dinámicos. La
relevancia del estudio reside en la transformación de datos crudos en
indicadores estadísticos (medidas de posición, dispersión y forma),
validando que el análisis exploratorio de datos (EDA) es una etapa
indispensable para garantizar la robustez de cualquier inferencia o
modelo predictivo posterior.

\section{Objetivos}

\subsection{Objetivo General}

Analizar estadísticamente el comportamiento de los precios de los
vehículos basándose en sus características técnicas y de uso, mediante
la aplicación de herramientas de estadística descriptiva y visualización
de datos en el lenguaje R.

\subsection{Objetivos Específicos}

\begin{enumerate}
  \item Caracterizar las variables cualitativas (Marca, Combustible, Transmisión) a través de tablas de frecuencias y sus graficos asociados.
  \item Determinar las medidas de tendencia central y dispersión para las variables cuantitativas (Precio, Kilometraje, Año).
  \item Identificar la presencia de valores atípicos (outliers) en los precios que puedan sesgar el análisis del mercado.
  \item Comparar la estructura de costos entre vehículos de combustión interna y tecnologías alternativas.
\end{enumerate}

\section{Cobertura}

\subsection{Horizontal}

La cobertura horizontal abarca las características de cada vehículo
registrado en la base de datos. Se consideran las dimensiones técnicas
del producto (marca, modelo, tipo de motor y transmisión) y las métricas
de valor y uso (precio de mercado, kilometraje y año de fabricación) que
definen su perfil comercial.

\subsection{Vertical}

La cobertura vertical del trabajo de investigación comprende un total de
2.500 observaciones (registros). Cada fila representa un vehículo único
disponible en el mercado de segunda mano. El alcance temporal de los
datos es longitudinal, abarcando modelos desde el año 2000 hasta el año
2024, lo que permite un análisis de la depreciación y evolución
tecnológica a corto y mediano plazo.

\begin{table}[ht]
\centering
\caption{Variables consideradas en el estudio} 
\resizebox{15cm}{!} {
\begin{tabular}{cll}
  \hline
  \textbf{TIPO} & \textbf{VARIABLE} & \textbf{DESCRIPCIÓN} \\ 
  \hline
  Categórica & Brand & Nombre del fabricante del carro \\ 
  Numérica & Year & Año en que fue fabricado el carro \\ 
  Numérica & Engine.Size & Cilindrad del motor del carro en litros  \\ 
  Categórica & Fuel.Type & Tipo de combustible o sistema de propulsion (e.g. Gasolina, Electrico) \\ 
  Categórica & Transmission & Tipo de transmision del carro (Manual, Automatica) \\ 
  Numérica & Mileage & Millas recorridas en el carro \\ 
  Categórica & Condition & Condicion del carro segun su uso (e.g. Nuevo, Usado) \\ 
  Numérica & Price & Precio de venta en dolares \\ 
  Categórica & Model & Nombre del modelo especifico del carro \\ 
   \hline
\end{tabular}
}
\end{table}

\section{Periodo de Referencia}

Los registros de precios de los carros presentes en la base de datos
usados son los precios de dicho carro al momento de realizar la
recolección de los datos, por lo que los precios son para el ano 2023.
Sin embargo, cabe acotar que los anos de producción de los carros
presentes en la lista van desde el 2000 al 2023. La base de datos
analizada actúa como una fotografía del mercado secundario actual,
reflejando cómo se valoran hoy los activos producidos en las últimas dos
décadas.

\begin{itemize}
  \item Permite capturar el ciclo de vida completo de un vehículo, desde su salida de concesionario hasta su etapa de obsolescencia.
  \item Facilita el estudio de la transición tecnológica, integrando modelos de combustión de principios de siglo con las innovaciones eléctricas e híbridas comercializadas en el último bienio. \end{itemize}

\section{Variables}

La variable principal en estudio (Variable Dependiente/Objetivo) es
Price, que representa el valor de mercado del vehículo medido en
unidades monetarias. Como variables explicativas o independientes se
analizan:

\begin{enumerate}
  \item Variables de Uso y Desgaste: \textbf{Year} y \textbf{Mileage}
  \item Variables Categóricas: \textbf{Brand}, \textbf{Model} y \textbf{Condition}
  \item Variables de Especificación Técnica: \textbf{Engine.Size}, \textbf{Fuel.Type} y \textbf{Transmission}
\end{enumerate}

\chapter{Marco Metódico}

En esta sección se presentan las metodologías, procedimientos
algorítmicos y técnicas de curaduría de datos implementadas para
transformar la muestra inicial en un conjunto de datos técnicamente
verosímil. Se detalla el tratamiento de la fuente de datos
precios\_carros.csv, los protocolos de limpieza (Data Wrangling) y el
procedimiento de imputación mediante fuentes externas utilizado para
corregir las inconsistencias detectadas en las variables de precio y
millaje. \newline \newline Igualmente se explican las tecnicas
estadisticas para posterior análisis utilizando el lenguaje de
programación R (versión 4.x) bajo el entorno de desarrollo RStudio,
empleando el paradigma ``Tidyverse'' para la manipulación y
transformación de estructuras de datos, así como funciones de
aleatorización controlada para el ajuste de la muestra.

\section{Bases metódicas utilizadas en el trabajo de investigación}

El diseño de la investigación es de carácter cuantitativo, no
experimental y de corte transversal con perspectiva histórica, dado que
se analizan datos observados en un punto del tiempo, pero que abarcan
modelos de distintos años de fabricación. A continuación, se describen
las etapas fundamentales del procesamiento:

\subsection{Procesamiento y Limpieza de Datos (ETL)}

La calidad de las conclusiones estadísticas depende directamente de la
pureza de la base de datos original. Tras una inspección inicial del
archivo precios\_carros.csv mediante el lenguaje R, se identificaron
anomalías estructurales y lógicas que requirieron una intervención
inmediata bajo los siguientes criterios: \newline \newline
\textbf{Depuración de Registros Incompletos}

En primera instancia, se procedió a la eliminación de 250 registros que
presentaban valores nulos (NA) en todos sus campos. Esta decisión se
tomó para evitar sesgos en el cálculo de los momentos estadísticos
(media, varianza) que podrían derivar en interpretaciones incorrectas de
los datos. \newline \newline \textbf{Tratamiento de Discrepancias
Lógicas y Técnicas}

Posteriormente, se realizó una auditoría de coherencia entre los
registros y las especificaciones técnicas reales de la industria
automotriz, aplicando los siguientes filtros de corrección:

\begin{itemize}

 \item Anacronismos de Comercialización: Se eliminaron o corrigieron registros donde el año de fabricación era previo al lanzamiento oficial del modelo (ej. modelos de Tesla registrados antes de 2015).

 \item Consistencia en el Sistema de Propulsion: Se detectaron errores de clasificación en el tipo de combustible. Se ajustaron aquellos vehículos listados como Eléctricos que poseían cilindrada (Engine Size) y viceversa, garantizando que los modelos eléctricos tuvieran un valor de NA en el tamaño del motor.

 \item Validación de Cilindrada: Se filtraron registros con tamaños de motor fuera de los rangos mecánicos posibles para el modelo específico (ej. motores de 8 litros en sedanes compactos), ajustándolos a los estándares del fabricante para mantener la representatividad de la muestra.
 \end{itemize}

\textbf{Reingeniería de Datos y Validación Externa}

Tras identificar que la base de datos original presentaba una
distribución de precios y kilometrajes excesivamente uniforme e
incoherente con el año de fabricación (ausencia de correlación entre
antigüedad y valor), se procedió a una fase de curaduría de datos para
aproximar la muestra a la realidad del mercado automotriz.

El proceso de ajuste se ejecutó bajo la siguiente metodología:

\begin{itemize}

 \item Recolección de Parámetros Reales: Se utilizó el portal especializado CarGurus como fuente de referencia externa para extraer métricas actualizadas de precios y millas recorridas según marca y modelo.

 \item Estratificación Temporal: Se segmentó la muestra en cohortes quinquenales (2000-2005, 2006-2010, 2011-2015, 2016-2020 y 2021-2023). Esta técnica de estratificación permitió capturar el comportamiento del valor del carro y su millaje en distintas etapas de su ciclo de vida.

 \item Muestreo de Referencia: Para cada estrato y modelo, se recolectó una submuestra representativa de aproximadamente 30 observaciones, estableciendo así los límites técnicos del mercado.

 \item Imputación Aleatoria Controlada: Utilizando el lenguaje de programación R, se reemplazaron los valores ilógicos originales mediante la generación de números aleatorios. Para garantizar la representatividad y evitar el ruido de los valores extremos, se restringió el rango de asignación al 75% central de los datos de millas y precios recopilados en la fase de investigación externa.
\end{itemize}

\subsection {Técnicas de Análisis Estadístico}

\begin{itemize}
 \item Análisis Univariado: Se utilizarán medidas de tendencia central (Media, Mediana) y de dispersión (Desviación Estándar, Coeficiente de Variación) para caracterizar el comportamiento individual del Precio, el Kilometraje y el Año. Se emplearán histogramas de frecuencias para observar la forma de la distribución (Simetría y Curtosis).

 \item Análisis Bivariado y Comparativo: Para cumplir con el objetivo de comparar marcas y tipos de combustible, se aplicará la técnica de Segmentación de Datos. Se utilizarán Diagramas de Caja (Boxplots), los cuales permiten visualizar la mediana, los cuartiles y, fundamentalmente, identificar la existencia de valores atípicos (outliers) en los precios de cada categoría.

 \item Análisis de Relación: Para evaluar la relación entre el uso (Kilometraje) y el valor (Precio), se emplearán Diagramas de Dispersión (Scatter Plots) y el cálculo del Coeficiente de Correlación de Pearson, determinando así la intensidad y dirección de la asociación entre ambas variables.
\end{itemize}

%%%%%%%%%%%%%%%%%%%%%%%%%%%%%%%%%%%%%%%%%%%%%%%%%%%%%
%                   REFERENCIAS                     %
%%%%%%%%%%%%%%%%%%%%%%%%%%%%%%%%%%%%%%%%%%%%%%%%%%%%%
% Mantenemos la estructura natbib de tu plantilla
\renewcommand{\bibname}{Lista de Referencias}
\bibliographystyle{apalike} 
\bibliography{bibliografia.bib} 

%%%%%%%%%%%%%%%%%%%%%%%%%%%%%%%%%%%%%%%%%%%%%%%%%%%%%
%                   APÉNDICES                       %
%%%%%%%%%%%%%%%%%%%%%%%%%%%%%%%%%%%%%%%%%%%%%%%%%%%%%

\end{document}
